\section{The non-GRS pencil and its quotient}\label{sec:pencil}

Let \(H_t\) be the ordered matrix in \((1.2)\), over a finite field of odd
order, and put
\[
G(t)=t^4-4t^3+7t^2-4t+1.
\]
The admitted non-GRS condition is
\[
t(t-1)(t^2-t+1)(t^2-3t+1)G(t)\ne0.                    \tag{4.1}
\]
The first four factors keep the displayed arc and quotient coordinates
nondegenerate; \(G=0\) is the conic, hence GRS, boundary.  Write
\(C_t=\ker H_t\) and \(\Psi_t=\Psi_{C_t}\), and let \(A,B,z\) be as in
\((1.3)\).

\begin{theorem}[Projective and Clifford quotient]\label{thm:lc-pencil}
For admitted \(t,u\), the following are equivalent:
\[
H_t\sim_{\mathrm{proj}}H_u,\qquad
C_t\sim_{\mathrm{mon}}C_u,\qquad
\Psi_t\sim_{\LC}\Psi_u,\qquad z(t)=z(u),
\]
allowing arbitrary party permutations.
\end{theorem}

\begin{proof}
Set
\[
 y(t)=\frac{(t-1)^2}{t}.
\]
Then \(A=-4t^2y\), \(B=t^2(y^2-1)\), and
\[
 z=\frac{(y-y^{-1})^2}{16}.
\]
It follows that \(z(t)=z(u)\) exactly when
\[
 y(u)\in\{y(t),-y(t),y(t)^{-1},-y(t)^{-1}\}.            \tag{4.2}
\]
Each of the four relations is realized by a direct projectivity and a
party permutation.  In order, one may use
\[
\begin{array}{c|c}
\text{relation}&\text{projectivity}\\ \hline
y(u)=y(t)&\diag(1,(1-u)/(1-t),u/t)\\
y(u)=-y(t)&\diag(1,(1-u)/(1-t),-u/t)\\
y(u)=y(t)^{-1}&
\begin{psmallmatrix}1&0&0\\0&0&(u-1)/t\\0&-u/(1-t)&0\end{psmallmatrix}\\
y(u)=-y(t)^{-1}&
\begin{psmallmatrix}1&0&0\\0&0&(1-u)/t\\0&-u/(1-t)&0\end{psmallmatrix}.
\end{array}
\]
The corresponding permutations are respectively the identity,
\((5\,6)\), \((3\,5)(4\,6)\), and \((3\,5\,4\,6)\).
All denominators and determinants are nonzero on the admitted locus.

Conversely, the ten signed products of complementary \(3\)-by-\(3\)
brackets of \(H_t\) form the multiset
\[
\{+A,+A,+A,-A,-A,-A,+B,+B,-B,-B\}.                    \tag{4.3}
\]
Projectivity and independent column rescaling multiply all ten products
by one common scalar; party permutation changes at most their common
sign.  The multiplicities three and two in \((4.3)\) recover
\((B/A)^2=z\).  Hence projective equivalence forces equality of \(z\).
The equivalence with monomial codes follows from
Section~\ref{sec:dictionary}, and projective equivalence supplies an LC
equivalence.

For the remaining direction, shorten the Pauli-label Lagrangian to every
four-set.  Projection between the local two-dimensional label planes
defines transition maps; comparing two supports gives a holonomy matrix.
The multiset of conjugacy-and-inversion invariants
\(\Tr(M)^2/\det M\) is preserved by LC and party permutation.  Its 450
entries have exact values
\[
\begin{array}{c|c}
4&90\\
-1/z&144\\
(2z+1)^2/z^2&24\\
\text{each root of }X^2-8X+(8-16z-1/z)&96.
\end{array}                                             \tag{4.4}
\]
Away from the GRS value \(z=-1/4\), the multiplicity pattern identifies
\(-1/z\), including every modular collision.  Thus LC equivalence forces
equality of \(z\).
\end{proof}

The map \(t\mapsto z\) has degree eight: the four transformations in
\((4.2)\), followed by the quadratic relation
\(t+t^{-1}=y+2\), account for its eight sheets.  The signed refinement
\(w=B/A\), \(z=w^2\), is not quantum data.  Both even and odd explicit
party permutations reverse \(w\), and their projectivities lift to
monomial local Cliffords.

\begin{proof}[Proof of Corollary~\ref{cor:lu-lc-pencil}]
Theorem~\ref{thm:lc-pencil} gives all equivalences except LU.  Every LC is
an LU, while Theorem~\ref{thm:lu-lc-rigidity} turns every LU intertwiner
between the states \(\Psi_t,\Psi_u\) into an LC intertwiner.
\end{proof}
